\section{Points a resoudre et problemes identifies}

\subsection{Points confirmes par le depot}

\begin{itemize}
  \item Plusieurs architectures coexistent: APEX, \path{voiture_system}, anciens chemins de simulation et code Python historique. Cette coexistence peut compliquer l'onboarding.
  \item Certaines documentations historiques ne decrivent plus exactement la pile APEX actuelle.
  \item Le fonctionnement reel depend de ports serie et de l'acces a \path{/sys/class/pwm}, donc d'une configuration materielle precise.
  \item Le debit LiDAR doit etre verifie selon le capteur utilise; des documents anciens mentionnent des valeurs differentes selon les voitures.
\end{itemize}

\subsection{Points deduits ou potentiels}

\begin{itemize}
  \item Le diagnostic \path{recognition_tour_diagnostic_report.md} decrit un run particulier avec un abandon de type perte de trajectoire et une possible saturation de direction. Ce constat est utile, mais ne suffit pas a conclure a un defaut general de la pile.
  \item La complexite de la pile de lancement peut rendre difficile l'identification rapide du noeud responsable lors d'un probleme terrain.
  \item Les differences entre simulation et vehicule reel peuvent provoquer des ecarts importants si les parametres capteurs et actionneurs ne sont pas recales regulierement.
\end{itemize}

\subsection{Element a completer par retour terrain}

Cette section reste a completer avec les problemes observes en pratique sur la Voiture Blue: symptomes exacts, conditions de test, version du code, parametres utilises, logs ROS, et decision prise apres analyse. Aucun inventaire terrain exhaustif et valide n'a ete trouve dans le depot.

